\documentclass[12pt]{article}
\usepackage[utf8]{inputenc}
\usepackage[T1]{fontenc}
\usepackage[margin=1in]{geometry}
\usepackage{setspace}
\usepackage{xcolor}
\usepackage[colorlinks=true, urlcolor=blue!60!black, linkcolor=blue!60!black]{hyperref}
\setlength{\parindent}{2em}
\setlength{\parskip}{1.25em}
\renewcommand{\baselinestretch}{1.0}
\begin{document}
\begin{center}
{\LARGE \textit{Shelby County v. Holder}, 570 U.S. 529 (2013)\par}
\vspace{0.8em}
{\large Case Brief --- Four-Page Edition\par}
\end{center}
\vspace{1em}

\section{Memory Jogger}

An Alabama county challenged the Voting Rights Act's preclearance regime; the Court held 5--4 that the \S~4(b) coverage formula, still based on decades-old data, could no longer justify subjecting some States to \S~5 preclearance, because current burdens must be justified by current needs and the equal sovereignty of the States (\href{https://www.law.cornell.edu/supremecourt/text/12-96}{Cornell LII}).

\section{Detailed Case Facts}

After nearly a century of southern nullification of the Fifteenth Amendment through literacy tests, poll taxes, and white primaries, Congress passed the Voting Rights Act of 1965. Section 2 permanently bars racially discriminatory voting practices nationwide. Section 4(b)'s formula covered jurisdictions that used a ``test or device'' in 1964 and had registration or turnout below 50\%. Section 5 barred covered jurisdictions from making \emph{any} voting change until the Attorney General or a three-judge D.C. court precleared it, with the jurisdiction bearing the burden of proof. \textit{South Carolina v. Katzenbach} (1966) upheld the scheme as ``rational in both practice and theory'' given ``exceptional conditions.''

Congress reauthorized the temporary provisions in 1970, 1975, 1982, and 2006, when, after a record exceeding 15,000 pages, it extended them 25 years by votes of 390--33 and 98--0 while leaving the 1964/1968/1972 formula untouched and strengthening \S~5. In \textit{Northwest Austin} (2009), the Court avoided the constitutional question but warned that ``[t]hings have changed in the South'' and that the Act's ``current burdens'' must be justified by ``current needs.''

\section{Procedural History}

Shelby County, Alabama, ineligible for bailout, sued Attorney General Eric Holder in 2010 for a declaration that \S\S~4(b) and 5 were facially unconstitutional. The district court granted summary judgment to the government. A divided D.C. Circuit affirmed, calling \S~4(b) ``a close question''; Judge Stephen Williams dissented, finding the formula ``irrational.'' The Supreme Court reversed on June 25, 2013.

\section{Judicial Votes}

Reversed, \textbf{5--4}.
\begin{itemize}
\item \textbf{Majority:} Chief Justice Roberts, joined by Justices Scalia, Kennedy, Thomas, and Alito.
\item \textbf{Concurrence:} Justice Thomas.
\item \textbf{Dissent:} Justice Ginsburg, joined by Justices Breyer, Sotomayor, and Kagan.
\end{itemize}

\section{Holding}

Section 4(b) is unconstitutional because its formula rests on 40-year-old data and eradicated practices that no longer speak to current conditions, so it cannot be used to subject jurisdictions to preclearance. The Court issued no holding on \S~5 and left \S~2 untouched, but without a formula, preclearance became inoperable unless Congress enacts a new one (apart from judicial ``bail-in'' under \S~3(c)).

\section{Analysis of Opinions}

\textbf{Majority (Roberts): federalism and equal sovereignty.} The Framers rejected a federal ``negative'' on state laws, and the States retain power to regulate elections. The ``fundamental principle of equal sovereignty'' ``remains highly pertinent in assessing subsequent disparate treatment of States.'' Covered States had to ``beseech the Federal Government for permission'' while neighbors acted freely, ``strong medicine'' that \textit{Katzenbach} tolerated only because of exceptional conditions.

\textbf{Changed conditions.} ``[N]early 50 years later, things have changed dramatically.'' Registration gaps had closed (Mississippi's from 63.2 points to $-3.8$), Black turnout exceeded white turnout in five of six originally covered States, and objection rates had fallen to 0.16\%. The deterrence argument was non-falsifiable.

\textbf{An irrational formula.} The formula's criteria themselves must be relevant; Congress could not ``reverse-engineer'' it. The Fifteenth Amendment ``is not designed to punish for the past.'' Congress did not use its record of ``second-generation barriers'' to craft a new formula, and those barriers only highlighted the mismatch.

\textbf{Thomas (concurring).} Section 5 itself is unconstitutional; leaving it intact ``needlessly prolongs the demise of that provision.''

\textbf{Ginsburg (dissenting).} The Reconstruction Amendments grant Congress affirmative power, at its zenith for voting. Discrimination evolved like ``battling the Hydra'' into second-generation barriers; DOJ blocked over 700 changes from 1982 to 2006, and covered jurisdictions accounted for 56\% of successful \S~2 suits. The majority entertained a facial challenge from a county with its own recent violations (``[h]ubris is a fit word''), silently overruled \textit{Katzenbach} on equal sovereignty, and ignored the record.

\section{Examples: Future Applications}

\begin{enumerate}
\item \textbf{Same-Side (Invalid): Poll-Tax History.} A formula covering States that used a poll tax before 1964 rests on history alone; ``history did not end in 1965.''
\item \textbf{Same-Side (Invalid): Low Turnout.} A formula covering States with turnout below 50\% uses fresh data but is not ``sufficiently related to the problem that it targets.''
\item \textbf{Opposite-Side (Likely Valid): Adjudicated Violations.} A rolling nationwide trigger based on five final judgments of intentional voting discrimination uses current data and treats States symmetrically.
\item \textbf{Opposite-Side (Valid): Judicial Bail-In.} After finding intentional discrimination, a court orders a county to preclear changes under \S~3(c), resting on a present-day finding.
\item \textbf{Fence-Sitter: Equal Sovereignty Elsewhere.} A sports-wagering license law favors two States. \textit{Shelby County} calls equal sovereignty ``highly pertinent,'' but its reach beyond preclearance is unsettled.
\end{enumerate}

\section{Critique}

\textbf{Text and doctrine.} ``Equal sovereignty'' appears nowhere in the Constitution's text, and the Reconstruction Amendments expressly altered antebellum federalism. The facial invalidation ignored \textit{Salerno} and the Act's severability clause despite Alabama's own record.

\textbf{Unresolved standard.} The Court left open whether Fifteenth Amendment legislation faces \textit{Katzenbach}'s rational-means test or \textit{Boerne}'s congruence and proportionality, an ambiguity \textit{Louisiana v. Callais} (2026) resolved in favor of \textit{Boerne}, holding that ``\S~2 imposes liability only when the evidence supports a strong inference'' of intentional discrimination.

\section{Key Quotations}

\begin{quote}
``Throwing out preclearance when it has worked and is continuing to work to stop discriminatory changes is like throwing away your umbrella in a rainstorm because you are not getting wet.''
\end{quote}
\noindent --- Justice Ginsburg, dissenting, 570 U.S. at 590.

\begin{quote}
``But history did not end in 1965. By the time the Act was reauthorized in 2006, there had been 40 more years of it. \ldots{} The Fifteenth Amendment commands that the right to vote shall not be denied or abridged on account of race or color, and it gives Congress the power to enforce that command. The Amendment is not designed to punish for the past; its purpose is to ensure a better future.''
\end{quote}
\noindent --- Chief Justice Roberts, for the Court, 570 U.S. at 552--53.

\begin{quote}
``The Voting Rights Act of 1965 employed extraordinary measures to address an extraordinary problem. Section 5 of the Act required States to obtain federal permission before enacting any law related to voting---a drastic departure from basic principles of federalism. And \S~4 of the Act applied that requirement only to some States---an equally dramatic departure from the principle that all States enjoy equal sovereignty.''
\end{quote}
\noindent --- Chief Justice Roberts, for the Court, 570 U.S. at 534--35.

\end{document}
